In this section $\mathbb S$ is again a $(\VF,\Gamma)$-generated ground with
$\Gamma_0=\dcl^{\mathrm{eq}}(\mathbb S)\cap\Gamma$ nontrivial, and
$g=[\G_m]\in R(\mathbb S)$. We consider definable families of objects of
$\VF_*$ indexed by a value-group parameter and compute the bounded value of
the corresponding inflated total space.

Let $n,m,s\geq0$ and let
\[
 \mathcal A\subseteq\Gamma^n\times K^m\times\RV_\infty^s
\]
be $\mathbb S$-definable, such that $\{t:(\gamma,a,t)\in\mathcal A\}$ is
finite for every $\gamma\in\Gamma^n$ and $a\in K^m$. Write
$A_\gamma=\{(a,t):(\gamma,a,t)\in\mathcal A\}$ for its rows. If
$\gamma\in\Gamma_0^n$, then $A_\gamma$ is an $\mathbb S$-definable object of
$\VF_*$. For arbitrary $\gamma\in\Gamma^n$, let $\mathbb S\langle\gamma\rangle$
be the ground obtained by adding the coordinates of $\gamma$ to the Gamma
parameters of $\mathbb S$; it is again $(\VF,\Gamma)$-generated, and
$A_\gamma$ is an $\mathbb S\langle\gamma\rangle$-definable object of
$\VF_*$. The inflated family is
\[
 A^{\mathrm{inf}}
   =\{(x,a,t)\in(K^\times)^n\times K^m\times\RV_\infty^s:\
       (\vv(x),a,t)\in\mathcal A\}.
\]
It is an $\mathbb S$-definable object of $\VF_*$: its RV fibre over $(x,a)$
is the RV fibre of $\mathcal A$ over $(\vv(x),a)$. For an RV family
$\mathcal U\subseteq\Gamma^n\times\RV_\infty^r$ we use the analogous RV
inflation
\[
 \mathcal U^{\mathrm{inf}}
   =\{(z,u)\in(\RV^\times)^n\times\RV_\infty^r:\ (\vrv(z),u)\in\mathcal U\}.
\]
The bounded value of an RV set is taken with any finite-to-one label, for
instance a coordinate projection; it does not depend on that choice.

\begin{thm}[Inflated-family Fubini]\label{thm:inflated}
There are a finite $\mathbb S$-definable partition $(E_a)_{a}$ of $\Gamma^n$
and classes $v_a\in R(\mathbb S)$ with the following properties.
\begin{enumerate}
\item If $A_\gamma=\varnothing$ for some $\gamma\in E_a$, then
$A_\gamma=\varnothing$ for every $\gamma\in E_a$, and $v_a=0$.
\item For every $a$ and every $\gamma\in E_a\cap\Gamma_0^n$,
\[
 \mathbb E_{\mathrm b}^{\mathbb S}\!\left(\int[A_\gamma]\right)=v_a .
\]
\item In $R(\mathbb S)$,
\begin{equation}\label{inflated:fubini}
 \mathbb E_{\mathrm b}^{\mathbb S}\!\left(\int[A^{\mathrm{inf}}]\right)
   =g^n\sum_a\chi_{\mathrm b}(E_a)\,v_a .
\end{equation}
\end{enumerate}
\end{thm}

The factor $g^n$ comes from the additional $n$ valued-field coordinates of
$A^{\mathrm{inf}}$. For $A_\gamma=X\cap\vv^{-1}(\gamma)$, with
$X\subseteq(K^\times)^n$, the theorem still concerns the inflated total
space rather than $X$ itself; compare the tropical motivic Fubini theorem
\cite[Theorem 3.1.3]{NicaisePayne}. No cancellation of $g^n$ is asserted.

The proof has four steps: uniform representatives over the value-group
parameter, a relative monomial decomposition, a separate partition of the
parameter space, and the final assembly.

\medskip
\noindent\emph{Uniform representatives.}

\begin{lem}\label{lem:inflated:uniform}
There are finitely many $\mathbb S$-definable RV families
$\mathcal U_l\subseteq\Gamma^n\times\RV_\infty^{r_l}$, $1\leq l\leq q$, and a
finite $\mathbb S$-definable partition $(P_l)_l$ of $\Gamma^n$ such that
$\mathcal U_l$ lies over $P_l$, and:
\begin{enumerate}
\item for every $\gamma\in\Gamma^n$, $A_\gamma=\varnothing$ if and only if
every row $\mathcal U_{l,\gamma}$ is empty;
\item for every $\gamma\in P_l\cap\Gamma_0^n$,
$\mathbb E_{\mathrm b}^{\mathbb S}(\int[A_\gamma])
 =\mathbb E_{\mathrm b}^{\mathbb S}([\mathcal U_{l,\gamma}])$, and
$\mathcal U_{l',\gamma}=\varnothing$ for $l'\neq l$;
\item $\mathbb E_{\mathrm b}^{\mathbb S}(\int[A^{\mathrm{inf}}])
 =\sum_l\mathbb E_{\mathrm b}^{\mathbb S}([\mathcal U_l^{\mathrm{inf}}])$.
\end{enumerate}
\end{lem}

\begin{proof}
Fix $\gamma\in\Gamma^n$. The lifting theorem over the generated ground
$\mathbb S\langle\gamma\rangle$ gives an
$\mathbb S\langle\gamma\rangle$-definable presentation in some degree
at most $k$ in the sense of Definition~\ref{defn:mixed:presentation},
whose underlying set $D_\gamma\subseteq\RV_\infty^{c+k}$ has its label as
the last $k$ coordinates, together with a bound for its label fibres and an
$\mathbb S\langle\gamma\rangle$-definable bijection
$H_\gamma:A_\gamma\to\mathbb L D_\gamma$ given by its graph.

A set definable over $\mathbb S\langle\gamma\rangle$ is the row at $\gamma$
of an $\mathbb S$-definable family over $\Gamma^n$: replace the new Gamma
parameters by variables, keeping the parameters of $\mathbb S$. Here the
family is a set of actual value-group tuples, defined through the full RV
preimage of the parameter, so no RV representative of $\gamma$ is chosen.
Apply this to the formulas for $D_\gamma$ and for the graph of $H_\gamma$.
For fixed formulas and a fixed bound, validity at a parameter $\gamma'$
asserts the bound for the label fibres of the row of the first family and
that the row of the second family is the graph of a bijection from
$A_{\gamma'}$ onto the lift of that row. It is a first-order condition on
$\gamma'$ over $\mathbb S$. These conditions cover $\Gamma^n$. Compactness
gives finitely many data, and a definable disjoint refinement of their
validity loci gives the partition $(P_l)_l$, with one datum on each $P_l$.
Let $\mathcal U_l$ be the restriction of its first family to $P_l$.

Property (1) holds because $H_\gamma$ is a bijection and a lift is empty
exactly when its underlying set is. For $\gamma\in P_l\cap\Gamma_0^n$ the
datum consists of $\mathbb S$-definable sets, and $H_\gamma$ is an
$\mathbb S$-definable bijection onto $\mathbb L\mathcal U_{l,\gamma}$, so (2)
follows from Lemma~\ref{lem:mixed:lift}. For (3), the uniform graphs give an
$\mathbb S$-definable bijection
\[
 A^{\mathrm{inf}}\cap\vv^{-1}(P_l)\ \xrightarrow{\ \sim\ }\ \mathbb L\,
 \mathcal U_l^{\mathrm{inf}},\qquad
 (x,a,t)\longmapsto\bigl(x,H_{\vv(x)}(a,t)\bigr),
\]
where $\mathcal U_l^{\mathrm{inf}}$ carries the label $(z,y)$, with $y$ the
label of the datum. Lemma~\ref{lem:mixed:lift} and finite additivity over
the partition $(\vv^{-1}(P_l))_l$ of $A^{\mathrm{inf}}$ prove (3).
\end{proof}

\medskip
\noindent\emph{A relative monomial decomposition.}

\begin{lem}\label{lem:relative:chart}
Let $X\subseteq(\RV^\times)^n\times(\RV^\times)^r$ be $\mathbb S$-definable,
with coordinates $(z,u)$, such that membership of $(z,u)$ in $X$ depends on
$z$ only through $\vrv(z)$. Then $X$ has a finite $\mathbb S$-definable
partition $(X_i)_i$ with the following data: $N_i\in\GL_r(\Z)$, integers
$h_i+\ell_i=r$, a nonempty $\mathbb S$-definable set
$W_i\subseteq(\RV^\times)^{h_i}$ with $\vrv(W_i)=\{\beta_i\}$ for some
$\beta_i\in\Gamma_0^{h_i}$, and an $\mathbb S$-definable set
$J_i\subseteq\Gamma^n\times\Gamma^{\ell_i}$, such that, writing
$h_{N_i}(u)=(w,y)$ with $w\in(\RV^\times)^{h_i}$,
\begin{equation}\label{relative:monomial:chart}
 (z,u)\in X_i\quad\Longleftrightarrow\quad
 w\in W_i\ \text{ and }\ (\vrv(z),\vrv(y))\in J_i .
\end{equation}
Moreover, if $(\gamma,\eta)\in J_i$, then $(z,u)\in X$ for some $(z,u)$ with
$\vrv(z)=\gamma$.
\end{lem}

\begin{proof}
Apply Lemma~\ref{coh:decom} to $X$ in $n+r$ coordinates. It gives value-induced
pieces $X'=X\cap\vrv^{-1}(D)$, matrices $M\in\GL_{n+r}(\Z)$ and
\[
 h_M(X')=V\times I^\sharp,\qquad V\subseteq(\RV^\times)^{m},\qquad
 \vrv(V)=\{\beta\}.
\]
Each piece is invariant under multiplication of $z$ by $(\mathbf k^\times)^n$,
because both $X$ and $D$ depend on $z$ only through its values. Let $B$ be the
integer matrix formed by the first $m$ rows and the first $n$ columns of $M$.
An element $c\in(\mathbf k^\times)^n$ acts on the first $m$ coordinates of
$h_M(z,u)$ by multiplication with $h_B(c)$.

The left kernel $\{y\in\Z^m:yB=0\}$ is a saturated subgroup of $\Z^m$, hence
a direct summand. Choose $Q\in\GL_m(\Z)$ whose first $h$ rows form a basis of
it. Then the first $h$ rows of $QB$ vanish and its remaining $p=m-h$ rows
have trivial left kernel. An integer matrix with trivial left kernel defines
a monomial map which is surjective on every divisible group; since $\mathbf k$
is algebraically closed, $\mathbf k^\times$ is divisible. Hence $(\mathbf
k^\times)^n$ acts on $h_Q(V)$ through the full group $\{1\}^h\times(\mathbf
k^\times)^p$. Since $h_Q(V)$ has one value tuple, invariance gives
\[
 h_Q(V)=W\times\beta_{\mathrm{free}}^\sharp ,
\]
where $W$ is the projection of $h_Q(V)$ to its first $h$ coordinates.

The first $h$ rows of $\operatorname{diag}(Q,1)M$ vanish in the $z$
columns. Their restrictions to the $u$ columns form an $h\times r$ integer
matrix $R$ with an integer right inverse, since
$\operatorname{diag}(Q,1)M\in\GL_{n+r}(\Z)$. Rows with an integer right
inverse extend to a matrix $N\in\GL_r(\Z)$ whose first $h$ rows are $R$.
Write $h_N(u)=(w,y)$. Then $w$ is the first $h$ coordinates of
$h_{\operatorname{diag}(Q,1)M}(z,u)$, and the values of all the other
coordinates of this tuple are integral linear functions of $\vrv(z)$ and
$\vrv(y)$, because $w$ has the fixed value tuple. Membership in the piece
is $w\in W$ together with a condition on the values of the remaining
coordinates. Therefore \eqref{relative:monomial:chart} holds with
\[
 J=\{(\vrv(z),\vrv(y)):(z,u)\in X'\},
\]
the value image of the piece. It is $\mathbb S$-definable, and no angular
representative is retained as a parameter. The set $W$ is a definable
projection, nonempty because $V$ is, and its value tuple is fixed over
$\mathbb S$. The last assertion holds because $J$ is a value image of a
subset of $X$.
\end{proof}

If $X\subseteq(\RV^\times)^n\times\RV_\infty^r$, apply the lemma to each
zero pattern of the last $r$ coordinates, after deleting the zero coordinates.
For an RV set as in the lemma and $\gamma\in\Gamma^n$, put
$X_\gamma=\{u:\exists z\ (\vrv(z)=\gamma,\ (z,u)\in X)\}$ and
$J_{i,\gamma}=\{\eta:(\gamma,\eta)\in J_i\}$. The coordinate change
$(z,u)\mapsto(z,h_{N_i}(u))$ maps $X_i$ onto $W_i\times J_i^\sharp$, after a
permutation of coordinates. The monomial evaluation
\eqref{eq:bounded-evaluation} therefore gives
\begin{equation}\label{relative:chart:values}
 \mathbb E_{\mathrm b}([X])
   =\sum_i[W_i]\,g^{n+\ell_i}\chi_{\mathrm b}(J_i),
 \qquad
 \mathbb E_{\mathrm b}([X_\gamma])
   =\sum_i[W_i]\,g^{\ell_i}\chi_{\mathrm b}(J_{i,\gamma})
 \quad(\gamma\in\Gamma_0^n),
\end{equation}
where the sums include the pieces of all zero patterns.

\medskip
\noindent\emph{A separate base partition.}
The projections of the sets $J_i$ to $\Gamma^n$ need not form a partition.
For example, in the constant family $\mathcal U_\gamma=\{1,2\rv(t),3\rv(t)\}$
over a ground containing $t$ with $\vv(t)\neq0$, the pieces at internal values
$0$ and $\vv(t)$ both project onto the whole parameter space. We therefore
construct the partition of the parameter space separately.

Let $J\subseteq\Gamma^n\times\Gamma^\ell$ be $\mathbb S$-definable and put
$e(\gamma)=\chi_{\mathrm b}(J_\gamma)$, the bounded Euler characteristic of
the slice. It is the integral over the last $\ell$ coordinates of the
indicator of $J$. By the Gamma fibre-integration formula of
Section~\ref{sec:setup}, together with its independence of the order of
coordinates, $e$ is an integer function on $\Gamma^n$ with finite range and
$\mathbb S$-definable level sets, and
\begin{equation}\label{gamma:projection:formula}
 \chi_{\mathrm b}(J)=\sum_{c}c\,\chi_{\mathrm b}(\{\gamma:e(\gamma)=c\}).
\end{equation}
If $J_\gamma=\varnothing$, then $e(\gamma)=0$.

\begin{proof}[Proof of Theorem~\ref{thm:inflated}]
Take $\mathcal U_l$ and $P_l$ from Lemma~\ref{lem:inflated:uniform}. Each
$X^{(l)}=\mathcal U_l^{\mathrm{inf}}$ satisfies the hypothesis of
Lemma~\ref{lem:relative:chart}, on every zero pattern of its last $r_l$
coordinates; its row at $\gamma$ is $\mathcal U_{l,\gamma}$. Index all the
resulting charts by pairs $(l,i)$, with data $W_{li}$, $\ell_{li}$ and
$J_{li}$. Let $(E_a)_a$ be the common finite refinement of the level sets of
all the functions $e_{li}(\gamma)=\chi_{\mathrm b}(J_{li,\gamma})$ and of the
sets $\{\gamma:\mathcal U_{l,\gamma}\neq\varnothing\}$. Write $e_{li,a}$ for
the value of $e_{li}$ on $E_a$, and put
\[
 c_{li}=[W_{li}]\,g^{\ell_{li}},\qquad
 v_a=\sum_{l,i}e_{li,a}\,c_{li}\qquad\text{in }R(\mathbb S).
\]

For (3), Lemma~\ref{lem:inflated:uniform}(3), the first formula of
\eqref{relative:chart:values} and \eqref{gamma:projection:formula} give
\[
 \mathbb E_{\mathrm b}\!\left(\int[A^{\mathrm{inf}}]\right)
 =\sum_{l,i}[W_{li}]\,g^{n+\ell_{li}}\chi_{\mathrm b}(J_{li})
 =g^n\sum_{l,i}c_{li}\sum_ae_{li,a}\chi_{\mathrm b}(E_a)
 =g^n\sum_a\chi_{\mathrm b}(E_a)\,v_a .
\]
For (2), let $\gamma\in E_a\cap\Gamma_0^n$. By
Lemma~\ref{lem:inflated:uniform}(2) and the second formula of
\eqref{relative:chart:values},
\[
 \mathbb E_{\mathrm b}\!\left(\int[A_\gamma]\right)
 =\sum_l\mathbb E_{\mathrm b}([\mathcal U_{l,\gamma}])
 =\sum_{l,i}c_{li}\,\chi_{\mathrm b}(J_{li,\gamma})
 =v_a .
\]
For (1), suppose that $A_{\gamma_0}=\varnothing$ for some $\gamma_0\in E_a$.
Then every row $\mathcal U_{l,\gamma_0}$ is empty. Since $E_a$ refines the
sets $\{\gamma:\mathcal U_{l,\gamma}\neq\varnothing\}$, every row
$\mathcal U_{l,\gamma}$ with $\gamma\in E_a$ is empty, and hence
$A_\gamma=\varnothing$. By the last assertion of
Lemma~\ref{lem:relative:chart}, each $J_{li}$ lies over
$\{\gamma:\mathcal U_{l,\gamma}\neq\varnothing\}$, so its slice at $\gamma_0$
is empty and $e_{li,a}=0$. Thus $v_a=0$.
\end{proof}
